\frfilename{mt264.tex}
\versiondate{12.5.03}
\copyrightdate{1994}

\def\chaptername{Perubahan variabel dalam integral}
\def\sectionname{Ukuran Hausdorff}

\newsection{264}

Pokok berikutnya yang ingin saya bahas adalah persoalan `ukuran
permukaan'; contoh yang berguna untuk terus diingat sepanjang bagian ini dan
bagian berikutnya ialah gagasan luas bagi daerah-daerah pada permukaan bola, tetapi
setiap permukaan dua dimensi lain yang melengkung mulus dalam ruang tiga
dimensi juga sama baiknya. Menurut saya, sangatlah masuk akal bahwa konsep
intuitif kita tentang `luas' permukaan demikian seharusnya bersesuaian dengan
ukuran yang tepat. Namun, memformalkan intuisi ini tidak sepele, terutama jika
kita menginginkan keumuman yang diarahkan oleh gagasan geometris sederhana;
maksud saya, kita tidak berpuas diri dengan argumen yang bergantung pada sifat
khusus permukaan bola, misalnya, untuk menguraikan luas permukaan bola. Saya membagi
persoalan ini menjadi dua bagian. Dalam bagian ini saya akan menguraikan suatu
konstruksi yang memungkinkan kita mendefinisikan ukuran berdimensi $r$ dari
suatu permukaan berdimensi $r$—di antara hal-hal lain—dalam ruang berdimensi
$s$. Dalam bagian berikutnya saya akan menyajikan teorema-teorema dasar yang
memungkinkan ukuran-ukuran ini dihitung secara efektif dalam kasus-kasus utama.

\leader{264A}{Definisi} Misalkan $s\ge 1$ suatu bilangan bulat, dan $r>0$.
\cmmnt{(Perhatian utama saya adalah $r$ yang bulat,
tetapi saya tidak akan mensyaratkan hal ini sebelum benar-benar diperlukan,
sebab ada beberapa gagasan yang sangat menarik yang melibatkan `dimensi'
tak bulat $r$.)} Untuk sebarang $A\subseteq \BbbR^s$, $\delta>0$, tetapkan

$$\eqalign{\theta_{r\delta}A
=\inf\{\sum_{n=0}^{\infty}(\diam A_n)^r:\sequencen{A_n}
&\text{ adalah barisan subhimpunan }\BbbR^s\text{ yang meliput }A,\cr
&\qquad\qquad\qquad\quad\diam A_n\le\delta\text{ untuk setiap }
  n\in\Bbb N\}.\cr}$$

\noindent Dalam konteks ini, akan memudahkan jika kita mengatakan bahwa
$\diam\emptyset=0$.
Sekarang tetapkan

\Centerline{$\theta_rA=\sup_{\delta>0}\theta_{r\delta}A$;}

\noindent $\theta_r$ adalah {\bf ukuran luar Hausdorff berdimensi $r$} pada
$\BbbR^s$.

\leader{264B}{}\cmmnt{ Tentu saja kita harus segera memeriksa hal berikut:

\medskip

\noindent}{\bf Lema} $\theta_r$, sebagaimana didefinisikan dalam 264A, selalu
merupakan ukuran luar.

\proof{ Anda seharusnya sudah terbiasa dengan argumen-argumen ini, tetapi
terdapat satu langkah tambahan dalam argumen ini, jadi saya menguraikan
rinciannya.

\medskip

{\bf (a)} Dengan menafsirkan diameter himpunan kosong sebagai $0$, kita
mempunyai $\theta_{r\delta}\emptyset=0$ untuk setiap $\delta>0$, sehingga
$\theta_r\emptyset=0$.

\medskip

{\bf (b)} Jika $A\subseteq B\subseteq\BbbR^s$, maka setiap barisan yang
meliput $B$ juga meliput $A$, sehingga
$\theta_{r\delta}A\le\theta_{r\delta}B$ untuk setiap $\delta$ dan
$\theta_rA\le\theta_rB$.

\medskip

{\bf (c)} Misalkan $\sequencen{A_n}$ suatu barisan subhimpunan $\BbbR^s$
yang gabungannya $A$, dan ambil sebarang $a<\theta_rA$. Maka terdapat
$\delta>0$ sedemikian sehingga $a\le\theta_{r\delta}A$. Sekarang
$\theta_{r\delta}A\le\sum_{n=0}^{\infty}\theta_{r\delta}(A_n)$. \Prf\
Ambil $\epsilon>0$, dan untuk setiap $n\in\Bbb N$ pilih suatu barisan
$\sequence{m}{A_{nm}}$ himpunan yang meliput $A_n$, dengan
$\diam A_{nm}\le\delta$ untuk setiap $m$ dan
$\sum_{m=0}^{\infty}(\diam A_{nm})^r\le\theta_{r\delta}A_n+2^{-n}\epsilon$.
Maka $\langle A_{nm}\rangle_{m,n\in\Bbb N}$ merupakan peliput $A$ oleh
terhitung banyak himpunan berdiameter paling besar $\delta$, sehingga

\Centerline{$\theta_{r\delta}A
\le\sum_{n=0}^{\infty}\sum_{m=0}^{\infty}(\diam A_{nm})^r
\le\sum_{n=0}^{\infty}(\theta_{r\delta}A_n+2^{-n}\epsilon)
=2\epsilon+\sum_{n=0}^{\infty}\theta_{r\delta}A_n$.}

\noindent Karena $\epsilon$ sebarang, kita memperoleh hasilnya.\ \Qed

Dengan demikian,

\Centerline{$a\le\theta_{r\delta}A
\le\sum_{n=0}^{\infty}\theta_{r\delta}A_n
\le\sum_{n=0}^{\infty}\theta_rA_n$.}

\noindent Karena $a$ sebarang,

\Centerline{$\theta_rA\le\sum_{n=0}^{\infty}\theta_rA_n$;}

\noindent karena $\sequencen{A_n}$ sebarang, $\theta_r$ merupakan ukuran
luar.
}%end of proof of 264B

\leader{264C}{Definisi} Jika $s\ge 1$ suatu bilangan bulat, dan $r>0$, maka
{\bf ukuran Hausdorff berdimensi $r$} pada $\BbbR^s$ adalah ukuran
$\mu_{Hr}$ pada $\BbbR^s$ yang didefinisikan melalui metode \Caratheodory\
dari ukuran luar $\theta_r$ dalam 264A-264B.

\leader{264D}{Catatan}\cmmnt{ {\bf (a)} Penting diperhatikan bahwa
himpunan-himpunan yang digunakan dalam definisi $\theta_{r\delta}$ tidak harus
berupa bola; bahkan dalam $\BbbR^2$, tidak setiap himpunan $A$ dapat diliput
oleh suatu bola yang berdiameter sama dengan $A$.

\header{264Db}}{\bf (b)} Dalam definisi-definisi di atas saya mensyaratkan
$r>0$. Kadang-kadang tepat untuk mengambil $\mu_{H0}$ sebagai ukuran cacah.
\cmmnt{Ini hampir merupakan hasil penerapan rumus-rumus di atas dengan $r=0$,
tetapi dapat timbul kesulitan jika kita menafsirkannya terlalu harfiah.}

\header{264Dc}{\bf (c)} Semua ukuran Hausdorff harus lengkap\cmmnt{, karena
semuanya didefinisikan melalui metode \Caratheodory\ (212A)}. Untuk $r>0$,
semuanya tanpa atom\cmmnt{ (264Yg)}. \cmmnt{Namun, ditinjau dari kriteria lain
dalam \S211, perilakunya sangat buruk; misalnya, jika $r$, $s$ bilangan bulat
dan $1\le r<s$, maka $\mu_{Hr}$ pada $\BbbR^s$ tidak semihingga. (Saya akan
memberikan pembuktiannya dalam 439H pada Jilid 4.) Meskipun demikian,
ukuran-ukuran tersebut memiliki beberapa sifat mencolok yang membuatnya cukup
mudah ditangani.}

\cmmnt{
\header{264Dd}{\bf (d)} Dalam 264A, perhatikan bahwa
$\theta_{r\delta}A\le\theta_{r\delta'}A$ ketika $0<\delta'\le\delta$;
akibatnya, misalnya,
$\theta_rA=\lim_{n\to\infty}\theta_{r,2^{-n}}A$. Saya telah mengizinkan
himpunan sebarang $A_n$ dalam peliput, tetapi hasilnya tidak berubah jika kita
membatasi perhatian pada peliput yang terdiri atas himpunan terbuka atau
himpunan tertutup (264Xc).
}%end of comment

\leader{264E}{Teorema} Misalkan $s\ge 1$ suatu bilangan bulat, dan $r\ge 0$;
misalkan $\mu_{Hr}$ ukuran Hausdorff berdimensi $r$ pada $\BbbR^s$, dan
$\Sigma_{Hr}$ domainnya. Maka setiap subhimpunan Borel dari $\BbbR^s$ termasuk
dalam $\Sigma_{Hr}$.

\proof{ Hal ini sepele jika $r=0$; jadi selanjutnya andaikan bahwa $r>0$.

\medskip

{\bf (a)} Langkah pertama ialah memperhatikan bahwa jika $A$, $B$ merupakan
subhimpunan $\Bbb R^s$ dan $\eta>0$ sedemikian sehingga
$\|x-y\|\ge\eta$ untuk semua $x\in A$, $y\in B$, maka
$\theta_r(A\cup B)=\theta_rA+\theta_rB$, dengan $\theta_r$ ukuran luar
Hausdorff berdimensi $r$ pada $\BbbR^s$. \Prf\ Tentu saja
$\theta_r(A\cup B)\le\theta_rA+\theta_rB$, karena $\theta_r$ merupakan ukuran
luar. Untuk ketaksamaan sebaliknya, kita dapat mengandaikan bahwa
$\theta_r(A\cup B)<\infty$, sehingga $\theta_rA$ dan $\theta_rB$ keduanya
berhingga. Ambil $\epsilon>0$ dan $\delta_1$, $\delta_2>0$ sedemikian sehingga

\Centerline{$\theta_rA+\theta_rB\le
\theta_{r\delta_1}A+\theta_{r\delta_2}B+\epsilon$.}

\noindent Tetapkan $\delta=\min(\delta_1,\delta_2,\bover12\eta)>0$ dan
misalkan $\sequencen{A_n}$ suatu barisan himpunan yang berdiameter paling
besar $\delta$, meliput $A\cup B$, dan sedemikian sehingga
$\sum_{n=0}^{\infty}(\diam A_n)^r\le\theta_{r\delta}(A\cup B)+\epsilon$.
Tetapkan

\Centerline{$K=\{n:A_n\cap A\ne\emptyset\}$,
\quad$L=\{n:A_n\cap B\ne\emptyset\}$.}

\noindent Karena

\Centerline{$\|x-y\|\ge\eta>\diam A_n$}

\noindent setiap kali $x\in A$, $y\in B$ dan $n\in\Bbb N$, kita mempunyai
$K\cap L=\emptyset$; dan tentu saja $A\subseteq\bigcup_{n\in K}A_n$,
$B\subseteq\bigcup_{n\in L}A_n$. Akibatnya,

$$\eqalign{\theta_rA+\theta_rB
&\le\epsilon+\theta_{r\delta_1}A+\theta_{r\delta_2}B\cr
&\le\epsilon+\sum_{n\in K}(\diam A_n)^r+\sum_{n\in L}(\diam A_n)^r\cr
&\le\epsilon+\sum_{n=0}^{\infty}(\diam A_n)^r
\le2\epsilon+\theta_{r\delta}(A\cup B)
\le 2\epsilon+\theta_r(A\cup B).\cr}$$

\noindent Karena $\epsilon$ sebarang, $\theta_r(A\cup
B)\ge\theta_rA+\theta_rB$, sebagaimana diperlukan.\ \Qed

\medskip

{\bf (b)} Diperoleh bahwa
$\theta_rA=\theta_r(A\cap G)+\theta_r(A\setminus G)$ setiap kali
$A\subseteq \BbbR^s$ dan $G$ terbuka.
\Prf\ Seperti biasa, cukup dipertimbangkan kasus $\theta_rA<\infty$ dan
ditunjukkan bahwa dalam kasus ini
$\theta_r(A\cap G)+\theta_r(A\setminus G)\le\theta_rA$. Tetapkan

\Centerline{$A_n=\{x:x\in A,\,\|x-y\|\ge 2^{-n}$ untuk setiap
$y\in A\setminus G\}$,}

\Centerline{$B_0=A_0$,
\quad$B_n=A_n\setminus A_{n-1}$ untuk $n>0$.}

\noindent Perhatikan bahwa $A_n\subseteq A_{n+1}$ untuk setiap $n$ dan
$\bigcup_{n\in\Bbb N}A_n=\bigcup_{n\in\Bbb N}B_n=A\cap G$. Intinya ialah
bahwa jika $m$, $n\in\Bbb N$ dan $n\ge m+2$, serta jika $x\in B_m$ dan
$y\in B_n$, maka terdapat $z\in A\setminus G$ sedemikian sehingga
$\|y-z\|<2^{-n+1}\le 2^{-m-1}$, sedangkan $\|x-z\|$ harus paling sedikit
$2^{-m}$, sehingga $\|x-y\|\ge\|x-z\|-\|y-z\|\ge 2^{-m-1}$. Akibatnya, untuk
setiap $k\ge 0$,

\Centerline{$\sum_{m=0}^k\theta_rB_{2m}
=\theta_r(\bigcup_{m\le k}B_{2m})\le\theta_r(A\cap G)<\infty$,}

\Centerline{$\sum_{m=0}^k\theta_rB_{2m+1}
=\theta_r(\bigcup_{m\le k}B_{2m+1})\le\theta_r(A\cap G)<\infty$,}

\noindent (dengan melakukan induksi pada $k$ dan menggunakan (a) di atas pada
langkah induksi). Akibatnya $\sum_{n=0}^{\infty}\theta_rB_n<\infty$.

Sekarang, jika $\epsilon>0$ diberikan, terdapat $m$ sedemikian sehingga
$\sum_{n=m}^{\infty}\theta_rB_n\le\epsilon$, sehingga

$$\eqalignno{\theta_r(A\cap G)+\theta_r(A\setminus G)
&\le\theta_rA_m+\sum_{n=m}^{\infty}\theta_rB_n+\theta_r(A\setminus G)\cr
&\le\epsilon+\theta_rA_m+\theta_r(A\setminus G)
=\epsilon+\theta_r(A_m\cup(A\setminus G))\cr
\noalign{\noindent (berdasarkan (a) lagi, karena $\|x-y\|\ge 2^{-m}$ untuk
$x\in A_m$, $y\in A\setminus G$)}
&\le\epsilon+\theta_rA.\cr}$$

\noindent Karena $\epsilon$ sebarang,
$\theta_r(A\cap G)+\theta_r(A\setminus G)\le\theta_rA$,
sebagaimana diperlukan.\ \Qed

\medskip

{\bf (c)} Bagian (b) menunjukkan tepat bahwa himpunan terbuka termasuk dalam
$\Sigma_{Hr}$. Langsung diperoleh bahwa aljabar-$\sigma$ Borel dari
$\BbbR^s$ termuat dalam $\Sigma_{Hr}$, sebagaimana dinyatakan.
}%end of proof of 264E

\leader{264F}{Proposisi} Misalkan $s\ge 1$ suatu bilangan bulat, dan $r>0$;
misalkan $\theta_r$ ukuran luar Hausdorff berdimensi $r$ pada $\BbbR^s$, dan
tuliskan $\mu_{Hr}$ untuk ukuran Hausdorff berdimensi $r$ pada $\BbbR^s$,
$\Sigma_{Hr}$ untuk domainnya. Maka

(a) untuk setiap $A\subseteq \BbbR^s$ terdapat suatu himpunan Borel
$E\supseteq A$ sedemikian sehingga $\mu_{Hr}E=\theta_rA$;

(b) $\theta_r=\mu^*_{Hr}$, yaitu ukuran luar yang didefinisikan dari
$\mu_{Hr}$;

(c) jika $E\in\Sigma_{Hr}$ dapat dinyatakan sebagai gabungan terhitung dari
himpunan-himpunan berukuran berhingga, terdapat himpunan Borel $E'$, $E''$
sedemikian sehingga $E'\subseteq E\subseteq E''$ dan
$\mu_{Hr}(E''\setminus E')=0$.

\proof{{\bf (a)} Jika $\theta_rA=\infty$, hal ini sepele—ambil
$E=\BbbR^s$. Jika tidak, untuk setiap $n\in\Bbb N$ pilih suatu barisan
$\sequence{m}{A_{nm}}$ himpunan berdiameter paling besar $2^{-n}$, yang
meliput $A$, dan sedemikian sehingga
$\sum_{m=0}^{\infty}(\diam A_{nm})^r\le\theta_{r,2^{-n}}A+2^{-n}$. Tetapkan
$F_{nm}=\overline{A}_{nm}$, $E=\bigcap_{n\in\Bbb N}\bigcup_{m\in\Bbb
N}F_{nm}$; maka $E$ merupakan suatu himpunan Borel dalam $\BbbR^s$. Tentu saja

\Centerline{$A\subseteq\bigcap_{n\in\Bbb N}\bigcup_{m\in\Bbb N}A_{nm}
\subseteq\bigcap_{n\in\Bbb N}\bigcup_{m\in\Bbb N}F_{nm}=E$.}

Untuk setiap $n\in\Bbb N$,

\Centerline{$\diam F_{nm}=\diam A_{nm}\le 2^{-n}$
untuk setiap $m\in\Bbb N$,}

\Centerline{$\sum_{m=0}^{\infty}(\diam F_{nm})^r
=\sum_{m=0}^{\infty}(\diam A_{nm})^r
\le\theta_{r,2^{-n}}A+2^{-n}$,}

\noindent sehingga

\Centerline{$\theta_{r,2^{-n}}E\le\theta_{r,2^{-n}}A+2^{-n}$.}

\noindent Dengan membiarkan $n\to\infty$,

\Centerline{$\theta_rE=\lim_{n\to\infty}\theta_{r,2^{-n}}E
\le\lim_{n\to\infty}\theta_{r,2^{-n}}A+2^{-n}
=\theta_rA$;}

\noindent tentu saja diperoleh bahwa $\theta_rA=\theta_rE$, karena
$A\subseteq E$. Sekarang, berdasarkan 264E kita mengetahui bahwa
$E\in\Sigma_{Hr}$, sehingga kita dapat menuliskan $\mu_{Hr}E$ sebagai
pengganti $\theta_rE$.

\medskip

{\bf (b)} Hal ini langsung diperoleh, karena kita mempunyai

\Centerline{$\mu_{Hr}^*A
=\inf\{\mu_{Hr}E:E\in\Sigma_{Hr},\,A\subseteq E\}
=\inf\{\theta_rE:E\in\Sigma_{Hr},\,A\subseteq E\}
\ge\theta_rA$}

\noindent untuk setiap $A\subseteq \BbbR^s$. Di pihak lain, jika
$A\subseteq\BbbR^s$, kita mempunyai suatu himpunan Borel $E\supseteq A$
sedemikian sehingga $\theta_rA=\mu_{Hr}E$, sehingga
$\mu_{Hr}^*A\le\mu_{Hr}E=\theta_rA$.

\medskip

{\bf (c)(i)} Andaikan terlebih dahulu bahwa $\mu_{Hr}E<\infty$. Berdasarkan
(a), terdapat himpunan Borel $E''\supseteq E$, $H\supseteq E''\setminus E$
sedemikian sehingga $\mu_{Hr}E''=\theta_rE$,

\Centerline{$\mu_{Hr}H=\theta_r(E''\setminus E)
=\mu_{Hr}(E''\setminus E)
=\mu_{Hr}E''-\mu_{Hr}E
=\mu_{Hr}E''-\theta_rE
=0$.}

\noindent Jadi, dengan menetapkan $E'=E''\setminus H$, kita memperoleh suatu
himpunan Borel yang termuat dalam $E$, dan

\Centerline{$\mu_{Hr}(E''\setminus E')\le\mu_{Hr}H=0$.}

\medskip

\quad{\bf (ii)} Untuk kasus umum, nyatakan $E$ sebagai $\bigcup_{n\in\Bbb
N}E_n$ dengan $\mu_{Hr}E_n<\infty$ untuk setiap $n$; ambil himpunan Borel
$E'_n$, $E_n''$ sedemikian sehingga $E_n'\subseteq E_n\subseteq E_n''$ dan
$\mu_{Hr}(E_n''\setminus E_n')=0$ untuk setiap $n$; dan tetapkan
$E'=\bigcup_{n\in\Bbb N}E'_n$, $E''=\bigcup_{n\in\Bbb N}E_n''$.
}%end of proof of 264F

\leader{264G}{Fungsi Lipschitz}\dvro{\bf: }{ Definisi ukuran Hausdorff tepat
disesuaikan dengan hasil berikut, yang bersesuaian dengan 262D.

\wheader{264G}{6}{2}{2}{48pt}

\noindent}{\bf Proposisi} Misalkan $m$, $s\ge 1$ bilangan bulat, dan
$\phi:D\to\BbbR^s$ suatu fungsi $\gamma$-Lipschitz, dengan $D$ suatu
subhimpunan $\BbbR^m$. Maka untuk sebarang $A\subseteq D$ dan $r\ge 0$,

\Centerline{$\mu_{Hr}^*(\phi[A])\le\gamma^r\mu_{Hr}^*A$}

\noindent untuk setiap $A\subseteq D$, dengan menuliskan $\mu_{Hr}^*$ untuk
ukuran luar Hausdorff berdimensi $r$ baik pada $\BbbR^m$ maupun pada
$\BbbR^s$.

\proof{{\bf (a)} Kasus $r=0$ sepele, sebab ketika itu $\gamma^r=1$ dan
$\mu_{Hr}^*A=\mu_{H0}A=\#(A)$ jika $A$ berhingga, dan $\infty$ jika tidak,
sedangkan $\#(\phi[A])\le\#(A)$.

\medskip

{\bf (b)} Jika $r>0$, ambil sebarang $\delta>0$. Tetapkan
$\eta=\delta/(1+\gamma)$ dan tinjau $\theta_{r\eta}:\Cal
P\BbbR^m\to[0,\infty]$, yang didefinisikan seperti dalam 264A. Berdasarkan
264Fb kita mengetahui bahwa

\Centerline{$\mu_{Hr}^*A=\theta_rA\ge\theta_{r\eta}A$,}

\noindent sehingga terdapat suatu barisan $\sequencen{A_n}$ himpunan,
semuanya berdiameter paling besar $\eta$, yang meliput $A$, dengan
$\sum_{n=0}^{\infty}(\diam A_n)^r\le\mu_{Hr}^*A+\delta$.
Sekarang $\phi[A]\subseteq\bigcup_{n\in\Bbb N}\phi[A_n\cap D]$ dan

\Centerline{$\diam \phi[A_n\cap D]
\le\gamma\diam A_n\le\gamma\eta\le\delta$}

\noindent untuk setiap $n$. Akibatnya,

\Centerline{$\theta_{r\delta}(\phi[A])\le\sum_{n=0}^{\infty}
(\diam \phi[A_n\cap D])^r
\le\sum_{n=0}^{\infty}\gamma^r(\diam A_n)^r
\le\gamma^r(\mu_{Hr}^*A+\delta)$,}

\noindent dan

\Centerline{$\mu_{Hr}^*(\phi[A])
=\lim_{\delta\downarrow 0}\theta_{r\delta}(\phi[A])
\le\gamma^r\mu_{Hr}^*A$,}

\noindent sebagaimana dinyatakan.
}%end of proof of 264G

\leader{264H}{}\cmmnt{ Langkah berikutnya ialah menghubungkan ukuran Hausdorff
berdimensi $r$ pada $\BbbR^r$ dengan ukuran Lebesgue pada $\BbbR^r$. Fakta
dasar yang kita perlukan adalah sebagai berikut, dan gagasan dalam
pembuktiannya bahkan lebih penting daripada hasilnya.

\wheader{264H}{6}{2}{2}{48pt}\noindent}{\bf Teorema}
%hold this line to get spacing right
Misalkan $r\ge 1$ suatu bilangan bulat, dan $A$ suatu subhimpunan terbatas
dari $\BbbR^r$; tuliskan $\mu_r$ untuk ukuran Lebesgue pada $\BbbR^r$ dan
$d=\diam A$. Maka

\Centerline{$\mu_r^*(A)\le\mu_rB(\tbf{0},\Bover{d}2)=2^{-r}\beta_rd^r$,}

\noindent dengan $B(\tbf{0},\bover{d}2)$ bola berpusat di $\tbf{0}$ dan
berdiameter $d$, sehingga $B(\tbf{0},1)$ merupakan bola satuan dalam
$\BbbR^r$, dan mempunyai ukuran

$$\eqalign{\beta_r
&=\Bover{1}{k!}\pi^k\text{ jika }r=2k\text{ genap},\cr
&=\Bover{2^{2k+1}k!}{(2k+1)!}\pi^k\text{ jika }r=2k+1\text{ ganjil}.\cr}$$

\proof{{\bf (a)} Untuk perhitungan $\beta_r$, lihat 252Q atau 252Xi.

\medskip

{\bf (b)} Kasus $r=1$ elementer, sebab dalam kasus ini $A$ termuat dalam suatu
interval dengan panjang $\diam A$, sehingga $\mu_1^*A\le\diam A$. Jadi,
selanjutnya andaikan bahwa $r\ge 2$.

\medskip

{\bf (c)} Untuk $1\le i\le r$, misalkan $S_i:\BbbR^r\to\BbbR^r$ refleksi
pada koordinat ke-$i$, sehingga
$S_ix=(\xi_1,\ldots,
\ifdim\pagewidth>467pt\penalty-100\fi
\xi_{i-1},-\xi_i,\xi_{i+1},
\ifdim\pagewidth>467pt\penalty-50\fi
\ldots,
\ifdim\pagewidth>467pt\penalty-10\fi
\xi_r)$ untuk setiap
$x=(\xi_1,\ldots,\xi_r)\in \BbbR^r$. Kita katakan bahwa suatu himpunan
$C\subseteq\BbbR^r$ {\bf simetris pada koordinat-koordinat dalam $J$}, dengan
$J\subseteq\{1,\ldots,r\}$, jika $S_i[C]=C$ untuk $i\in J$. Sekarang, inti
argumennya adalah fakta berikut: jika $C\subseteq\Bbb R^r$ suatu himpunan
terbatas yang simetris pada koordinat-koordinat dalam $J$, dengan $J$ suatu
subhimpunan sejati dari $\{1,\ldots,r\}$, dan
$j\in\{1,\ldots,r\}\setminus J$, maka terdapat suatu himpunan $D$, yang
simetris pada koordinat-koordinat dalam $J\cup\{j\}$, sedemikian sehingga
$\diam D\le\diam C$ dan $\mu_r^*C\le\mu^*_rD$.


\Prf\ {\bf (i)} Karena ukuran Lebesgue invarian terhadap permutasi koordinat,
cukup ditangani kasus $j=r$. Mulailah dengan menuliskan
$F=\overline{C}$, sehingga $\diam F=\diam C$ dan
$\mu_rF\ge\mu_r^*C$. Perhatikan bahwa karena $S_i$ merupakan homeomorfisme
untuk setiap $i$,

\Centerline{$S_i[F]=S_i[\overline{C}]=\overline{S_i[C]}=\overline{C}=F$}

\noindent untuk $i\in J$, dan $F$ simetris pada koordinat-koordinat dalam
$J$.

Untuk $y=(\eta_1,\ldots,\eta_{r-1})\in\BbbR^{r-1}$, tetapkan

\Centerline{$F_y=\{\xi:(\eta_1,\ldots,\eta_{r-1},\xi)\in F\}$,\quad
$f(y)=\mu_1F_y$,}

\noindent dengan $\mu_1$ ukuran Lebesgue pada $\Bbb R$.
Tetapkan

\Centerline{$D=\{(y,\xi):y\in\BbbR^{r-1},\,|\xi|<\Bover12f(y)\}
\subseteq\BbbR^r$.}

\medskip

\quad{\bf (ii)} Jika $H\subseteq\BbbR^r$ terukur dan $H\supseteq D$, maka,
dengan menuliskan $\mu_{r-1}$ untuk ukuran Lebesgue pada $\BbbR^{r-1}$, kita
mempunyai

$$\eqalignno{\mu_rH
&=\int\mu_1\{\xi:(y,\xi)\in H\}\mu_{r-1}(dy)\cr
\noalign{\noindent (dengan menggunakan 251N dan 252D)}
&\ge\int\mu_1\{\xi:(y,\xi)\in D\}\mu_{r-1}(dy)
=\int f(y)\mu_{r-1}(dy)\cr
&=\int\mu_1\{\xi:(y,\xi)\in F\}\mu_{r-1}(dy)
=\mu_rF
\ge\mu^*_rC.\cr}$$

\noindent Karena $H$ sebarang, $\mu_r^*D\ge\mu_r^*C$.

\medskip

\quad{\bf (iii)} Langkah berikutnya ialah memeriksa bahwa
$\diam D\le\diam C$. Jika $x$, $x'\in D$, nyatakan keduanya sebagai
$(y,\xi_r)$ dan $(y',\xi_r')$. Karena $F$ suatu himpunan tertutup terbatas
dalam $\BbbR^r$, $F_y$ dan $F_{y'}$ merupakan subhimpunan tertutup terbatas
dari $\Bbb R$. Selain itu, baik $f(y)$ maupun $f(y')$ harus lebih besar dari
$0$, sehingga $F_y$, $F_{y'}$ keduanya tak kosong. Akibatnya,

\Centerline{$\alpha=\inf F_y$,\quad $\beta=\sup F_y$,
\quad$\alpha'=\inf F_{y'}$,
\quad$\beta'=\sup F_{y'}$}

\noindent semuanya terdefinisi dalam $\Bbb R$, dan $\alpha$,
$\beta\in F_y$, sedangkan $\alpha'$ dan $\beta'$ termasuk dalam $F_{y'}$.
Kita mempunyai

$$\eqalign{|\xi_r-\xi_r'|
&\le|\xi_r|+|\xi'_r|
<\Bover12f(y)+\Bover12f(y')\cr
&=\Bover12(\mu_1F_y+\mu_1F_{y'})
\le\Bover12(\beta-\alpha+\beta'-\alpha')
\le\max(\beta'-\alpha,\beta-\alpha').\cr}$$

\noindent Jadi, dengan mengambil $(\xi,\xi')$ sebagai salah satu dari
$(\alpha,\beta')$ atau $(\beta,\alpha')$, kita dapat menemukan
$\xi\in F_y$, $\xi'\in F_{y'}$ sedemikian sehingga
$|\xi-\xi'|\ge|\xi_r-\xi'_r|$. Sekarang $z=(y,\xi)$, $z'=(y',\xi')$ keduanya
termasuk dalam $F$, sehingga

\Centerline{$\|x-x'\|^2=\|y-y'\|^2+|\xi_r-\xi_r'|^2
\le\|y-y'\|^2+|\xi-\xi'|^2
=\|z-z'\|^2
\le(\diam F)^2$,}

\noindent dan $\|x-x'\|\le\diam F$. Karena $x$ dan $x'$ sebarang,
$\diam D\le\diam F=\diam C$, sebagaimana dinyatakan.

\medskip

\quad{\bf (iv)} Jelas bahwa $S_r[D]=D$. Selain itu, jika $i\in J$, maka
(dengan menafsirkan $S_i$ sebagai operator pada $\BbbR^{r-1}$)

\Centerline{$F_{S_i(y)}=F_y$ untuk setiap $y\in\BbbR^{r-1}$,}

\noindent sehingga $f(S_i(y))=f(y)$ dan, untuk $\xi\in \Bbb R$,
$y\in\BbbR^{r-1}$,

\Centerline{$(y,\xi)\in D\,\iff\,|\xi|<\bover12 f(y)\,
\iff\,|\xi|<\bover12f(S_i(y))
\,\iff\,(S_i(y),\xi)\in D$,}

\noindent sehingga $S_i[D]=D$. Jadi $D$ simetris pada koordinat-koordinat
dalam $J\cup\{r\}$.\ \Qed

\medskip

{\bf (d)} Sisanya mudah. Mulai dengan sebarang himpunan terbatas
$A\subseteq\BbbR^r$; tetapkan $A_0=A$ dan secara induktif konstruksikan
$A_1,\ldots,A_r$ sedemikian sehingga

\Centerline{$d=\diam A=\diam A_0\ge\diam A_1\ge\ldots\ge\diam A_r$,}

\Centerline{$\mu_r^*A=\mu_r^*A_0\le\ldots\le\mu_r^*A_r$,}

\Centerline{$A_j$ simetris pada koordinat-koordinat dalam
$\{1,\ldots,j\}$ untuk setiap $j\le r$.}

\noindent Pada akhirnya, kita mempunyai $A_r$ yang simetris pada
koordinat-koordinat dalam $\{1,\ldots,r\}$. Namun, ini berarti bahwa jika
$x\in A_r$, maka

\Centerline{$-x=S_1S_2\ldots S_rx\in A_r$,}

\noindent sehingga

\Centerline{$\|x\|=\Bover12\|x-(-x)\|
\le\Bover12\diam A_r\le\Bover{d}2$.}

\noindent Jadi $A_r\subseteq B(\tbf{0},\bover{d}2)$, dan

\Centerline{$\mu_r^*A\le\mu_r^*A_r\le\mu_rB(\tbf{0},\Bover{d}2)$,}

\noindent sebagaimana dinyatakan.
}%end of proof of 264H


\leader{264I}{Teorema} Misalkan $r\ge 1$ suatu bilangan bulat; misalkan
$\mu$ ukuran Lebesgue pada $\BbbR^r$, dan misalkan $\mu_{Hr}$ ukuran Hausdorff
berdimensi $r$ pada $\BbbR^r$. Maka $\mu$ dan $\mu_{Hr}$ mempunyai
himpunan-himpunan terukur yang sama dan

\Centerline{$\mu E=2^{-r}\beta_r\mu_{Hr}E$}

\noindent untuk setiap himpunan terukur $E\subseteq\BbbR^r$, dengan
$\beta_r=\mu B(\tbf{0},1)$, sehingga faktor normalisasinya adalah

$$\eqalign{2^{-r}\beta_r
&=\Bover{1}{2^{2k}k!}\pi^k\text{ jika }r=2k\text{ genap},\cr
&=\Bover{k!}{(2k+1)!}\pi^k\text{ jika }r=2k+1\text{ ganjil}.\cr}$$

\proof{{\bf (a)} Tentu saja jika $B=B(x,\alpha)$ sebarang bola berjari-jari
$\alpha$,

\Centerline{$2^{-r}\beta_r(\diam B)^r
=\beta_r\alpha^r=\mu B$.}

\medskip

{\bf (b)} Intinya ialah bahwa $\mu^*=2^{-r}\beta_r\mu_{Hr}^*$. \Prf\
Misalkan $A\subseteq\BbbR^r$.

\medskip

\quad{\bf (i)} Misalkan $\delta$, $\epsilon>0$. Berdasarkan 261F, terdapat
suatu barisan $\sequencen{B_n}$ bola, semuanya berdiameter paling besar
$\delta$, sedemikian sehingga $A\subseteq\bigcup_{n\in\Bbb N}B_n$ dan
$\sum_{n=0}^{\infty}\mu B_n\le\mu^*A+\epsilon$. Sekarang, dengan
mendefinisikan $\theta_{r\delta}$ seperti dalam 264A,

\Centerline{$2^{-r}\beta_r\theta_{r\delta}(A)
\le 2^{-r}\beta_r\sum_{n=0}^{\infty}(\diam B_n)^r
=\sum_{n=0}^{\infty}\mu B_n
\le\mu^*A+\epsilon$.}

\noindent Dengan membiarkan $\delta\downarrow 0$,

\Centerline{$2^{-r}\beta_r\mu^*_{Hr}A\le\mu^*A+\epsilon$.}

\noindent Karena $\epsilon$ sebarang,
$2^{-r}\beta_r\mu_{Hr}^*A\le\mu^*A$.

\medskip

\quad{\bf (ii)} Misalkan $\epsilon>0$. Maka terdapat suatu barisan
$\sequencen{A_n}$ himpunan berdiameter paling besar $1$ sedemikian sehingga
$A\subseteq\bigcup_{n\in\Bbb N}A_n$ dan
$\sum_{n=0}^{\infty}(\diam A_n)^r\le\theta_{r1}A+\epsilon$, sehingga

\Centerline{$\mu^*A\le\sum_{n=0}^{\infty}\mu^*A_n
\le\sum_{n=0}^{\infty}2^{-r}\beta_r(\diam A_n)^r
\le 2^{-r}\beta_r(\theta_{r1}A+\epsilon)
\le 2^{-r}\beta_r(\mu_{Hr}^*A+\epsilon)$}

\noindent berdasarkan 264H. Karena $\epsilon$ sebarang,
$\mu^*A\le 2^{-r}\beta_r\mu_{Hr}^*A$.\ \Qed

\medskip

{\bf (c)} Karena $\mu$, $\mu_{Hr}$ merupakan ukuran yang didefinisikan dari
ukuran luarnya masing-masing melalui metode \Caratheodory, langsung diperoleh
bahwa $\mu=2^{-r}\beta_r\mu_{Hr}$ dalam pengertian ketat yang diperlukan.
}%end of proof of 264I

\leader{*264J}{Himpunan Cantor \dvrocolon{}}\cmmnt{ Dalam 264A saya menyatakan
bahwa `dimensi' pecahan $r$ menarik. Saya tidak mempunyai ruang untuk
membahasnya di sini, dan pokok itu berada di luar jalur utama jilid ini,
tetapi saya akan memberikan satu hasil karena hasil itu menarik pada dirinya
sendiri.

\medskip

\noindent}{\bf Proposisi} Misalkan $C$ himpunan Cantor dalam $[0,1]$.
Tetapkan $r=\ln 2/\ln 3$. Maka ukuran Hausdorff berdimensi $r$ dari $C$ adalah
$1$.

\proof{{\bf (a)} Ingat bahwa $C=\bigcap_{n\in\Bbb N}C_n$, dengan setiap
$C_n$ terdiri atas $2^n$ interval tertutup yang panjangnya $3^{-n}$, dan
$C_{n+1}$ diperoleh dari $C_n$ dengan menghapus sepertiga tengah (terbuka) dari
setiap interval dalam $C_n$. (Lihat 134G.) Karena $C$ tertutup,
$\mu_{Hr}C$ terdefinisi (264E). Perhatikan bahwa $3^r=2$.

\medskip

{\bf (b)} Jika $\delta>0$, ambil $n$ sedemikian sehingga
$3^{-n}\le\delta$; maka $C$ dapat diliput oleh $2^n$ interval berdiameter
$3^{-n}$, sehingga

\Centerline{$\theta_{r\delta}C\le 2^n(3^{-n})^r=1$.}

\noindent Akibatnya,

\Centerline{$\mu_{Hr}C=\mu_{Hr}^*C
=\lim_{\delta\downarrow 0}\theta_{r\delta}C\le 1$.}

\medskip

{\bf (c)} Kita memerlukan fakta elementer berikut: jika $\alpha$, $\beta$,
$\gamma\ge 0$ dan $\max(\alpha,\gamma)\le\beta$, maka
$\alpha^r+\gamma^r\le(\alpha+\beta+\gamma)^r$. \Prf\ Karena
$0<r\le 1$,

\Centerline{$\xi\mapsto(\xi+\eta)^r-\xi^r
=r\biggerint_0^{\eta}(\xi+\zeta)^{r-1}d\zeta$}

\noindent tak menaik untuk setiap $\eta\ge 0$.
Akibatnya,

$$\eqalign{(\alpha+\beta+\gamma)^r-\alpha^r-\gamma^r
&\ge (\beta+\beta+\gamma)^r-\beta^r-\gamma^r\cr
&\ge(\beta+\beta+\beta)^r-\beta^r-\beta^r
=\beta^r(3^r-2)
=0,\cr}$$

\noindent sebagaimana diperlukan.\ \Qed

\medskip

{\bf (d)} Sekarang andaikan bahwa $I\subseteq\Bbb R$ sebarang interval, dan
$m\in\Bbb N$; tuliskan $j_m(I)$ untuk banyaknya interval penyusun $C_m$ yang
termuat dalam $I$. Maka $2^{-m}j_m(I)\le(\diam I)^r$.
\Prf\ Jika $I$ tidak bertemu $C_m$, hal ini sepele. Jika tidak, lakukan induksi
pada

\Centerline{$l=\min\{i:I$ hanya bertemu dengan satu interval penyusun
$C_{m-i}\}$.}

\noindent Jika $l=0$, sehingga $I$ hanya bertemu dengan satu interval penyusun
$C_m$, maka $j_m(I)\le 1$, dan jika $j_m(I)=1$, maka
$\diam I\ge 3^{-m}$ sehingga $(\diam I)^r\ge 2^{-m}$; dengan demikian induksi
dimulai. Untuk langkah induksi menuju $l>0$, misalkan $J$ interval dalam
$C_{m-l}$ yang bertemu dengan $I$, dan $J'$, $J''$ kedua interval dalam
$C_{m-l+1}$ yang termuat dalam $J$, sehingga $I$ bertemu dengan $J'$ dan
$J''$, dan


\Centerline{$j_m(I)=j_m(I\cap J)=j_m(I\cap J')+j_m(I\cap J'')$.}

\noindent Berdasarkan hipotesis induksi,

\Centerline{$(\diam(I\cap J'))^r+(\diam(I\cap J''))^r
\ge 2^{-m}j_m(I\cap J')+2^{-m}j_m(I\cap J'')
=2^{-m}j_m(I)$.}

\noindent Di pihak lain, berdasarkan (c),

$$\eqalign{(\diam(I\cap J'))^r+(\diam(I\cap J''))^r
&\le(\diam(I\cap J')+3^{-m+l-1}+\diam(I\cap J''))^r\cr
&=(\diam(I\cap J))^r
\le(\diam I)^r\cr}$$

\noindent karena $J'$, $J''$ keduanya berdiameter paling besar
$3^{-(m-l+1)}$, yaitu panjang interval di antara keduanya.
Dengan demikian induksi berlanjut.\ \Qed

\medskip

{\bf (e)} Sekarang andaikan bahwa $\epsilon>0$. Maka terdapat suatu barisan
$\sequencen{A_n}$ himpunan yang meliput $C$, sedemikian sehingga
\discrcenter{468pt}{$\sum_{n=0}^{\infty}(\diam A_n)^r
<\mu_{Hr}C+\epsilon$.  }Ambil $\eta_n>0$ sedemikian sehingga
$\sum_{n=0}^{\infty}(\diam A_n+\eta_n)^r\le\mu_{Hr}C+\epsilon$, dan untuk
setiap $n$ ambil suatu interval terbuka $I_n\supseteq A_n$ dengan panjang
paling besar $\diam A_n+\eta_n$ dan dengan kedua titik ujung tidak termasuk
dalam $C$; hal ini mungkin karena $C$ tidak memuat interval nontrivial mana
pun. Sekarang $C\subseteq\bigcup_{n\in\Bbb N}I_n$; karena $C$ kompak,
terdapat $k\in\Bbb N$ sedemikian sehingga
$C\subseteq\bigcup_{n\le k}I_n$. Berikutnya, terdapat $m\in\Bbb N$ sedemikian
sehingga tidak ada titik ujung $I_n$, untuk $n\le k$, yang termasuk dalam
$C_m$. Akibatnya, setiap interval penyusun $C_m$ harus termuat dalam suatu
$I_n$, dan (dengan terminologi (d) di atas)
$\sum_{n=0}^kj_m(I_n)\ge 2^m$. Dengan demikian,

\Centerline{$1
\le\sum_{n=0}^k2^{-m}j_m(I_n)
\le \sum_{n=0}^k(\diam I_n)^r
\le \sum_{n=0}^{\infty}(\diam A_n+\eta_n)^r
\le\mu_{Hr}C+\epsilon$.}

\noindent Karena $\epsilon$ sebarang, $\mu_{Hr}C\ge 1$, sebagaimana
diperlukan.
}%end of proof of 264J

\leader{*264K}{Ruang metrik umum}\cmmnt{ Walaupun bab ini secara eksklusif
membahas ruang Euklides, pembaca yang mengenal teori umum ruang metrik mungkin
akan melihat sifat teori ini dengan lebih jelas jika mereka menggunakan
bahasa ruang metrik dalam definisi dan hasil dasar. Karena itu, saya mengulang
definisinya di sini dan menguraikan hasil-hasil yang bersesuaian dalam latihan
264Yb-264Yl.

} Misalkan $(X,\rho)$ suatu ruang metrik, dan $r>0$. Untuk sebarang
$A\subseteq X$, $\delta>0$, tetapkan

$$\eqalign{\theta_{r\delta}A
=\inf\{\sum_{n=0}^{\infty}(\diam A_n)^r:\sequencen{A_n}
&\text{ adalah barisan subhimpunan }X\text{ yang meliput }A,\cr
&\qquad\qquad\qquad\quad\diam A_n\le\delta\text{ untuk setiap }
  n\in\Bbb N\},\cr}$$

\noindent dengan menafsirkan diameter himpunan kosong sebagai $0$, dan
$\inf\emptyset$ sebagai $\infty$, sehingga $\theta_{r\delta}A=\infty$ jika
$A$ tidak dapat diliput oleh suatu barisan himpunan berdiameter paling besar
$\delta$. Katakan bahwa $\theta_rA=\sup_{\delta>0}\theta_{r\delta}A$ adalah
{\bf ukuran luar Hausdorff berdimensi $r$} dari $A$, dan ambil ukuran
$\mu_{Hr}$ yang didefinisikan melalui metode \Caratheodory\ dari ukuran luar
ini sebagai {\bf ukuran Hausdorff berdimensi $r$} pada $X$.

\exercises{
\leader{264X}{Latihan dasar $\pmb{>}$(a)}
%\spheader 264Xa
Tunjukkan bahwa semua fungsi $\theta_{r\delta}$ dalam 264A merupakan ukuran
luar. Tunjukkan bahwa dalam konteks tersebut,
$\theta_{r\delta}(A)=0$ jika dan hanya jika $\theta_r(A)=0$, untuk sebarang
$\delta>0$ dan sebarang $A\subseteq\BbbR^s$.

\spheader 264Xb Misalkan $s\ge 1$ suatu bilangan bulat, dan $\theta$ suatu
ukuran luar pada $\BbbR^s$ sedemikian sehingga
$\theta(A\cup B)=\theta A+\theta B$ setiap kali $A$, $B$ merupakan
subhimpunan tak kosong dari $\BbbR^s$ dan
$\inf_{x\in A,y\in B}\|x-y\|>0$. Tunjukkan bahwa setiap subhimpunan Borel dari
$\BbbR^s$ diukur oleh ukuran yang didefinisikan dari $\theta$ melalui metode
\Caratheodory.

\sqheader 264Xc Misalkan $s\ge 1$ suatu bilangan bulat dan $r>0$;
definisikan $\theta_{r\delta}$ seperti dalam 264A. Tunjukkan bahwa untuk
sebarang $A\subseteq\BbbR^s$, $\delta>0$,

$$\eqalign{\theta_{r\delta}A
&=\inf\{\sum_{n=0}^{\infty}(\diam F_n)^r:\sequencen{F_n}
 \text{ adalah barisan subhimpunan tertutup dari }\BbbR^s\cr
&\qquad\qquad\qquad\qquad\qquad\qquad\text{ yang meliput }A,\diam
F_n\le\delta
  \text{ untuk setiap }n\in\Bbb N\}\cr
&=\inf\{\sum_{n=0}^{\infty}(\diam G_n)^r:\sequencen{G_n}
  \text{ adalah barisan subhimpunan terbuka dari }\BbbR^s\cr
&\qquad\qquad\qquad\qquad\qquad\qquad\text{ yang meliput }A,\diam
G_n\le\delta
  \text{ untuk setiap }n\in\Bbb N\}.\cr}$$

\sqheader 264Xd Misalkan $s\ge 1$ suatu bilangan bulat dan $r\ge 0$;
misalkan $\mu_{Hr}$ ukuran Hausdorff berdimensi $r$ pada $\BbbR^s$.
Tunjukkan bahwa untuk setiap $A\subseteq \BbbR^s$ terdapat suatu himpunan
G$_{\delta}$ (yakni, suatu himpunan yang dapat dinyatakan sebagai irisan suatu
barisan himpunan terbuka) $H\supseteq A$ sedemikian sehingga
$\mu_{Hr}H=\mu^*_{Hr}A$. ({\it Petunjuk\/}: gunakan 264Xc.)

\sqheader 264Xe Misalkan $s\ge 1$ suatu bilangan bulat, dan $0\le r<r'$.
Tunjukkan bahwa jika $A\subseteq\BbbR^s$ dan ukuran luar Hausdorff berdimensi
$r$, yaitu $\mu_{Hr}^*A$, dari $A$ berhingga, maka $\mu_{Hr'}^*A$ harus nol.

\spheader 264Xf(i) Andaikan bahwa $f:[a,b]\to\Bbb R$ mempunyai grafik
$\Gamma_f\subseteq\BbbR^2$, dengan $a\le b$ dalam $\Bbb R$. Tunjukkan bahwa
ukuran luar $\mu_{H1}^*(\Gamma_f)$ dari $\Gamma_f$ untuk ukuran Hausdorff
berdimensi satu pada $\BbbR^2$ paling besar
$b-a+\Var_{[a,b]}(f)$.
\Hint{jika $f$ mempunyai variasi berhingga, tunjukkan bahwa
$\diam(\Gamma_{f\restr\ooint{t,u}})
\le u-t+\Var_{\ooint{t,u}}(f)$; kemudian gunakan 224E.} (ii) Misalkan
$f:[0,1]\to[0,1]$ fungsi Cantor (134H). Tunjukkan bahwa
$\mu_{H1}(\Gamma_f)=2$. \Hint{264G.}
%264G

\spheader 264Xg\dvAnew{2009.} Dalam 264A, tunjukkan bahwa

$$\eqalign{\theta_{r\delta}A
=\inf\{\sum_{n=0}^{\infty}(\diam A_n)^r:\sequencen{A_n}
&\text{ adalah barisan himpunan konveks yang meliput }A,\cr
&\qquad\qquad\qquad\quad\diam A_n\le\delta\text{ untuk setiap }
  n\in\Bbb N\}\cr}$$

\noindent untuk sebarang $A\subseteq\BbbR^s$.
%264A

\leader{264Y}{Latihan lanjutan (a)}
%\spheader 264Ya
Misalkan $\theta_{11}$ ukuran luar pada $\BbbR^2$ yang didefinisikan dalam
264A, dengan $r=\delta=1$, dan $\mu_{11}$ ukuran yang diturunkan dari
$\theta_{11}$ melalui metode \Caratheodory, serta $\Sigma_{11}$ domainnya.
Tunjukkan bahwa setiap himpunan dalam $\Sigma_{11}$ terabaikan atau
koterabaikan.

\spheader 264Yb Misalkan $(X,\rho)$ suatu ruang metrik dan $r\ge 0$.
Tunjukkan bahwa jika $A\subseteq X$ dan $\mu_{Hr}^*A<\infty$, maka $A$
separabel.

\spheader 264Yc Misalkan $(X,\rho)$ suatu ruang metrik, dan $\theta$ suatu
ukuran luar pada $X$ sedemikian sehingga
$\theta(A\cup B)=\theta A+\theta B$ setiap kali $A$, $B$ merupakan subhimpunan
tak kosong dari $X$ dan $\inf_{x\in A,y\in B}\rho(x,y)>0$. (Ukuran luar
demikian disebut {\bf ukuran luar metrik}.) Tunjukkan bahwa setiap subhimpunan
terbuka dari $X$ diukur oleh ukuran yang didefinisikan dari $\theta$ melalui
metode \Caratheodory.

\spheader 264Yd Misalkan $(X,\rho)$ suatu ruang metrik dan $r>0$;
definisikan $\theta_{r\delta}$ seperti dalam 264K. Tunjukkan bahwa untuk
sebarang $A\subseteq X$,

$$\eqalign{\mu_{Hr}^*A
&=\sup_{\delta>0}\inf\{\sum_{n=0}^{\infty}(\diam F_n)^r:\sequencen{F_n}
  \text{ adalah barisan subhimpunan tertutup dari }X\cr
&\qquad\qquad\qquad\qquad\qquad\qquad\text{ yang meliput }A,\,\diam
F_n\le\delta
  \text{ untuk setiap }n\in\Bbb N\}\cr
&=\sup_{\delta>0}\inf\{\sum_{n=0}^{\infty}(\diam G_n)^r:\sequencen{G_n}
  \text{ adalah barisan subhimpunan terbuka dari }X\cr
&\qquad\qquad\qquad\qquad\qquad\qquad\text{ yang meliput }A,\,\diam
G_n\le\delta
  \text{ untuk setiap }n\in\Bbb N\}.\cr}$$
\spheader 264Ye Misalkan $(X,\rho)$ suatu ruang metrik dan $r\ge 0$;
misalkan $\mu_{Hr}$ ukuran Hausdorff berdimensi $r$ pada $X$.
Tunjukkan bahwa untuk setiap $A\subseteq X$ terdapat suatu himpunan
G$_{\delta}$ $H\supseteq A$ sedemikian sehingga
$\mu_{Hr}H=\mu^*_{Hr}A$ merupakan ukuran luar Hausdorff berdimensi $r$ dari
$A$.

\spheader 264Yf Misalkan $(X,\rho)$ suatu ruang metrik dan $r\ge 0$;
misalkan $Y$ sebarang subhimpunan $X$, dan lengkapi $Y$ dengan metrik
terinduksi $\rho_Y$. (i) Tunjukkan bahwa ukuran luar Hausdorff berdimensi $r$
$\mu^{(Y)*}_{Hr}$ pada $Y$ tepat merupakan pembatasan pada $\Cal PY$ dari
ukuran luar $\mu_{Hr}^*$ pada $X$. (ii) Tunjukkan bahwa jika {\it salah satu}
dari $\mu_{Hr}^*Y<\infty$ {\it atau} $\mu_{Hr}$ mengukur $Y$ berlaku, maka
ukuran Hausdorff berdimensi $r$ $\mu^{(Y)}_{Hr}$ pada $Y$ tepat merupakan
ukuran subruang pada $Y$ yang diinduksi oleh ukuran $\mu_{Hr}$ pada $X$.

\spheader 264Yg Misalkan $(X,\rho)$ suatu ruang metrik dan $r>0$.
Tunjukkan bahwa ukuran Hausdorff berdimensi $r$ pada $X$ bersifat tanpa atom.
({\it Petunjuk\/}: Misalkan $E\in\dom\mu_{Hr}$. (i) Jika $E$ tidak separabel,
terdapat suatu himpunan terbuka $G$ sedemikian sehingga $E\cap G$ dan
$E\setminus G$ keduanya tidak separabel, sehingga keduanya tidak terabaikan.
(ii) Jika terdapat $x\in E$ sedemikian sehingga
$\mu_{Hr}(E\cap B(x,\delta))>0$ untuk setiap $\delta>0$, maka salah satu dari
himpunan-himpunan ini mempunyai komplemen yang tidak terabaikan dalam $E$.
(iii) Jika tidak, $\mu_{Hr}E=0$.)

\spheader 264Yh Misalkan $(X,\rho)$ suatu ruang metrik dan $r\ge 0$;
misalkan $\mu_{Hr}$ ukuran Hausdorff berdimensi $r$ pada $X$. Tunjukkan bahwa
jika $\mu_{Hr}E<\infty$, maka
$\mu_{Hr}E=\sup\{\mu_{Hr}F:F\subseteq E$ tertutup dan terbatas total$\}$.
({\it Petunjuk\/}: jika $\epsilon>0$ diberikan, gunakan 264Yd untuk menemukan
suatu himpunan tertutup dan terbatas total $F$ sedemikian sehingga
$\mu_{Hr}(F\setminus E)=0$ dan $\mu_{Hr}(E\setminus F)\le\epsilon$, lalu
terapkan 264Ye pada $F\setminus E$.)

\spheader 264Yi Misalkan $(X,\rho)$ suatu ruang metrik lengkap dan $r\ge 0$;
misalkan $\mu_{Hr}$ ukuran Hausdorff berdimensi $r$ pada $X$. Tunjukkan bahwa
jika $\mu_{Hr}E<\infty$, maka
$\mu_{Hr}E=\sup\{\mu_{Hr}F:F\subseteq E$ kompak$\}$.

\spheader 264Yj Misalkan $(X,\rho)$ dan $(Y,\sigma)$ ruang metrik. Jika
$D\subseteq X$ dan $\phi:D\to Y$ suatu fungsi, maka $\phi$ adalah
{\bf $\gamma$-Lipschitz} jika
$\sigma(\phi(x),\phi(x'))\le\gamma\rho(x,x')$ untuk setiap $x$, $x'\in D$.
(i) Tunjukkan bahwa dalam hal ini, jika $r\ge 0$,
$\mu_{Hr}^*(\phi[A])\le\gamma^r\mu_{Hr}^*A$ untuk setiap $A\subseteq D$,
dengan menuliskan $\mu_{Hr}^*$ untuk ukuran luar Hausdorff berdimensi $r$
pada $X$ maupun $Y$. (ii) Tunjukkan bahwa jika $X$ lengkap dan
$E\subseteq D$ dan $\mu_{Hr}E$ terdefinisi dan berhingga, maka
$\mu_{Hr}(\phi[E])$ terdefinisi.
({\it Petunjuk\/}: 264Yi.)

\spheader 264Yk Misalkan $(X,\rho)$ suatu ruang metrik, dan untuk $r\ge 0$
misalkan $\mu_{Hr}$ ukuran Hausdorff berdimensi $r$ pada $X$. Tunjukkan bahwa
terdapat tepat satu $\Delta=\Delta(X)\in[0,\infty]$ sedemikian sehingga
$\mu_{Hr}X=\infty$ jika $r\in\coint{0,\Delta}$, dan $0$ jika
$r\in\ooint{\Delta,\infty}$.

\spheader 264Yl Misalkan $(X,\rho)$ suatu ruang metrik dan $\phi:I\to X$
suatu fungsi kontinu, dengan $I\subseteq\Bbb R$ suatu interval. Tuliskan
$\mu_{H1}$ untuk ukuran Hausdorff berdimensi satu pada $X$. Tunjukkan bahwa

\Centerline{$\mu_{H1}(\phi[I])
\le\sup\{\sum_{i=1}^n\rho(\phi(t_i),\phi(t_{i-1})):t_0,\ldots,t_n\in
I,\,t_0\le\ldots\le t_n\}$,}

\noindent yaitu panjang kurva $\phi$, dengan kesamaan jika $\phi$ injektif.


\spheader 264Ym Tetapkan $r=\ln 2/\ln 3$, seperti dalam 264J, dan tuliskan
$\mu_{Hr}$ untuk ukuran Hausdorff berdimensi $r$ pada himpunan Cantor $C$.
Misalkan $\lambda$ ukuran biasa pada $\{0,1\}^{\Bbb N}$ (254J).
Definisikan $\phi:\{0,1\}^{\Bbb N}\to C$ dengan menetapkan
$\phi(x)=\bover23\sum_{n=0}^{\infty}3^{-n}x(n)$ untuk
$x\in\{0,1\}^{\Bbb N}$. Tunjukkan bahwa $\phi$ merupakan isomorfisme antara
$(\{0,1\}^{\Bbb N},\lambda)$ dan $(C,\mu_{Hr})$, sehingga $\mu_{Hr}$ merupakan
ukuran subruang pada $C$ yang diinduksi oleh `ukuran Cantor' sebagaimana
didefinisikan dalam 256Hc.

\spheader 264Yn Tetapkan $r=\ln 2/\ln 3$ dan tuliskan $\mu_{Hr}$ untuk ukuran
Hausdorff berdimensi $r$ pada himpunan Cantor $C$. Misalkan
$f:[0,1]\to[0,1]$ fungsi Cantor dan misalkan $\mu$ ukuran Lebesgue pada
$\Bbb R$. Tunjukkan bahwa $\mu f[E]=\mu_{Hr}E$ untuk setiap
$E\in\dom\mu_{Hr}$ dan $\mu_{Hr}(C\cap f^{-1}[F])=\mu F$ untuk setiap
himpunan terukur Lebesgue $F\subseteq[0,1]$.

\spheader 264Yo Misalkan $(X,\rho)$ suatu ruang metrik dan
$h:\coint{0,\infty}\to\coint{0,\infty},\ h(0)=0$ suatu fungsi tak menurun. Untuk
$A\subseteq X$, tetapkan

$$\eqalign{\theta_hA
&=\sup_{\delta>0}\inf\{\sum_{n=0}^{\infty}h(\diam A_n):\sequencen{A_n}
  \text{ adalah barisan subhimpunan dari }X\cr
&\qquad\qquad\qquad\qquad\qquad\qquad\text{ yang meliput }A,\,\diam
A_n\le\delta
  \text{ untuk setiap }n\in\Bbb N\},\cr}$$

\noindent dengan menafsirkan $\diam\emptyset$ sebagai $0$ dan
$\inf\emptyset$ sebagai $\infty$ seperti biasa. Tunjukkan bahwa $\theta_h$
merupakan ukuran luar pada $X$. Nyatakan dan buktikan teorema-teorema yang
bersesuaian dengan 264E dan 264F. Periksa 264X dan 264Y untuk mencari hasil
lain yang mungkin dapat digeneralisasi, barangkali dengan asumsi bahwa $h$
kontinu dari kanan.

\spheader 264Yp Misalkan $(X,\rho)$ suatu ruang metrik. Kita katakan bahwa
jika $a<b$ dalam $\Bbb R$ dan $f:[a,b]\to X$ suatu fungsi, maka $f$
{\bf kontinu mutlak} jika untuk setiap $\epsilon>0$ terdapat $\delta>0$
sedemikian sehingga $\sum_{i=0}^n\rho(f(a_i),f(b_i))\le\epsilon$ setiap kali
$a\le a_0\le b_0\le\ldots\le a_n\le b_n\le b$ dan
$\sum_{i=0}^n(b_i-a_i)\le\delta$. Tunjukkan bahwa $f:[a,b]\to X$ kontinu
mutlak jika dan hanya jika fungsi itu kontinu dan bervariasi terbatas (dalam
pengertian 224Ye) serta $\mu_{H1}f[A]=0$ setiap kali $A\subseteq[a,b]$
terabaikan Lebesgue, dengan $\mu_{H1}$ ukuran Hausdorff berdimensi satu pada
$X$. (Bandingkan 225M.)
Tunjukkan bahwa dalam hal ini $\mu_{H1}f[\,[a,b]\,]<\infty$.
%-  mt26bits

\spheader 264Yq Misalkan $s\ge 1$ suatu bilangan bulat, dan
$r\in\coint{s,\infty}$. Untuk $x$, $y\in\BbbR^s$, tetapkan
$\rho(x,y)=\|x-y\|^{s/r}$. (i) Tunjukkan bahwa $\rho$ merupakan metrik pada
$\BbbR^s$ yang menginduksi topologi Euklides. (ii) Misalkan $\mu_{Hr}$ ukuran
Hausdorff berdimensi $r$ yang bersesuaian. Tunjukkan bahwa
$\mu_{Hr}B(\tbf{0},1)=2^s$.
%264I out of order query
}%end of exercises

\endnotes{
\Notesheader{264} Dalam bagian ini kita telah sampai pada langkah berikutnya
dalam `teori ukuran geometris'. Saya membahasnya dengan sangat perlahan,
karena terdapat kesulitan nyata dalam pokok ini, dan untuk keperluan jilid ini
kita tidak perlu menguasai banyak darinya. Gagasannya di sini ialah menemukan
suatu definisi ukuran Lebesgue berdimensi $r$ yang `geometris' dalam
pengertian ketat, yakni hanya bergantung pada struktur metrik $\BbbR^r$, dan
karena itu dapat diterapkan pada himpunan yang mempunyai struktur metrik
tetapi tidak mempunyai struktur linear. Seperti yang sudah terjadi
sebelumnya, definisi ukuran Hausdorff dari suatu ukuran luar tidak
menimbulkan masalah—satu-satunya gagasan baru dalam 264A-264C ialah
penggunaan supremum $\theta_r=\sup_{\delta>0}\theta_{r\delta}$ dari
ukuran-ukuran luar—dan bagian sulitnya adalah membuktikan bahwa ukuran baru
kita mempunyai sifat yang berguna. Mengenai sifat ukuran Hausdorff, terdapat
dua sasaran pokok: pertama, memeriksa bahwa secara umum ukuran-ukuran ini
memiliki bagian yang cukup besar dari sifat ukuran Lebesgue; dan kedua,
membenarkan istilah `ukuran berdimensi $r$' dengan menghubungkan ukuran
Hausdorff berdimensi $r$ pada $\BbbR^r$ dengan ukuran Lebesgue pada
$\BbbR^r$.

Mengenai sifat ukuran Hausdorff umum, kita harus melangkah cukup hati-hati.
Saya tidak memberikan contoh tandingan di sini karena contoh-contoh itu
melibatkan konsep yang termasuk dalam Jilid 4 dan 5, bukan jilid ini, tetapi
saya harus memperingatkan Anda agar mengantisipasi hal terburuk. Setidaknya,
kita mengetahui bahwa himpunan terbuka terukur, sehingga semua himpunan Borel
terukur (264E). Ukuran luar suatu himpunan $A$ dapat didefinisikan melalui
himpunan-himpunan Borel yang memuat $A$ (264Fa), walaupun secara umum tidak
melalui himpunan-himpunan terbuka yang memuat $A$; tetapi ukuran suatu
himpunan terukur $E$ tidak selalu merupakan supremum ukuran
himpunan-himpunan Borel yang termuat dalam $E$, kecuali $E$ berukuran
berhingga (264Fc). Kita memang mendapati bahwa ukuran luar $\theta_r$ yang
didefinisikan dalam 264A merupakan ukuran luar yang didefinisikan dari
$\mu_{Hr}$ (264Fb), sehingga frasa `ukuran luar Hausdorff berdimensi $r$'
tidak ambigu. Salah satu sifat penting ukuran Lebesgue ialah bahwa ukuran
suatu himpunan terukur $E$ merupakan supremum ukuran subhimpunan kompak dari
$E$; sifat ini secara umum tidak dimiliki ukuran Hausdorff, tetapi berlaku
untuk himpunan $E$ berukuran berhingga dalam ruang lengkap (264Yi). Mengenai
subruang, tidak terdapat masalah dengan ukuran luar, dan untuk himpunan
berukuran berhingga, ukuran subruangnya juga konsisten (264Yf). Karena ukuran
Hausdorff didefinisikan dalam istilah metrik, ia berperilaku teratur di bawah
pemetaan Lipschitz (264G); salah satu kelas fungsi yang paling alami untuk
dipertimbangkan ketika mempelajari ruang metrik adalah fungsi $1$-Lipschitz,
sehingga (dalam bahasa 264G) $\mu_{Hr}^*\phi[A]\le \mu_{Hr}^*A$ untuk setiap
$A$.

Sifat pokok kedua ukuran Hausdorff, yakni hubungannya dengan ukuran Lebesgue
dalam dimensi yang sesuai, adalah Teorema 264I. Karena ukuran Hausdorff dan
ukuran Lebesgue sama-sama invarian terhadap translasi, hal ini dapat
dibuktikan dengan cara yang relatif elementer, kecuali dalam evaluasi
konstanta normalisasi; kita hanya perlu mengetahui bahwa
$\mu\coint{0,1}^r=1$ dan $\mu_{Hr}\coint{0,1}^r$ keduanya berhingga dan tidak
nol, dan ini langsung. (Argumen bagian (a) pembuktian 261F relevan.) Untuk
keperluan bab ini, menurut saya kita tidak harus mengetahui nilai konstanta
tersebut; tetapi saya tidak dapat membiarkannya tak terselesaikan, dan karena
itu memberikan Teorema 264H, yaitu {\bf ketaksamaan isodiametrik}, untuk
menunjukkan bahwa konstanta itu tepat merupakan ukuran Lebesgue dari bola
berdimensi $r$ dengan diameter $1$, sebagaimana diharapkan. Langkah kritis
dalam argumen 264H terdapat pada bagian (c) pembuktiannya. Langkah ini disebut
`simetrisasi Steiner'; gagasannya ialah bahwa, setelah diberikan suatu
himpunan $A$, kita mentransformasi $A$ melalui serangkaian langkah, pada
setiap tahap menurunkan, atau setidaknya tidak menaikkan, diameternya, dan
menaikkan, atau setidaknya tidak menurunkan, ukuran luarnya, dengan secara
bertahap membuat $A$ lebih simetris, sampai akhirnya diperoleh suatu himpunan
yang cukup terkendala sehingga mudah ditangani. Operasi simetrisasi khusus
yang digunakan dalam pembuktian ini cukup penting; tetapi gagasan regularisasi
bertahap terhadap suatu objek merupakan salah satu metode paling ampuh dalam
teori ukuran, dan Anda patut mencurahkan seluruh perhatian untuk menguasai
setiap contoh yang Anda jumpai. Berdasarkan pengalaman saya, gagasan itu
terutama berguna ketika mencari suatu ketaksamaan yang melibatkan besaran
yang berbeda sifat—dalam contoh sekarang, diameter dan volume suatu himpunan.

Tentu saja tidak nyaman memiliki dua ukuran pada $\BbbR^r$ yang berbeda
sebesar suatu kelipatan konstan, dan untuk keperluan bagian berikutnya
sebenarnya akan sedikit lebih mudah mengikuti {\smc Federer 69} dengan
menggunakan `ukuran Hausdorff ternormalisasi'
$2^{-r}\beta_r\mu_{Hr}$. (Untuk $r$ tak bulat, kita dapat mengambil
$\beta_r=\pi^{r/2}/\Gamma(1+\bover{r}2)$, sebagaimana disarankan dalam
252Xi.) Namun, saya yakin ini merupakan pandangan minoritas, dan contoh
mencolok ukuran Hausdorff pada himpunan Cantor (264J, 264Ym-264Yn) tampak jauh
lebih baik dalam versi yang tidak dinormalisasi.

Ukuran Hausdorff berdimensi ($\ln 2/\ln 3$) pada himpunan Cantor tentu
hanyalah satu, barangkali yang termudah, dari suatu kelas contoh yang besar.
Karena ukuran luar Hausdorff berdimensi $r$ dari suatu himpunan $A$, jika
dipandang sebagai fungsi $r$, berperilaku secara drastis (turun dari $\infty$
ke $0$) pada suatu nilai kritis $\Delta(A)$ (lihat 264Xe, 264Yk), ukuran itu
memberi kita suatu invarian ruang metrik dari $A$; $\Delta(A)$ adalah
{\bf dimensi Hausdorff} dari $A$. Jelas bahwa dimensi Hausdorff dari $C$
adalah $\ln 2/\ln 3$, sedangkan dimensi Hausdorff ruang Euklides berdimensi
$r$ adalah $r$.
}%end of notes

\discrpage
